Inverse aerodynamic design, given a target surface pressure or isentropic
Mach distribution, recovering the geometry that produces it, is classically
posed either as a well-posed but restrictive conformal-mapping problem
\citep{lighthill1945new,volpe1986design}, or as an outer optimisation loop
wrapped around a converged forward flow solve. The latter is the practical
default in industry: a geometry parameterisation, a surrogate or
gradient-based optimiser, and hundreds to thousands of fully converged flow
solves per design.

The motivation for the present paper is autobiographical. The author's own
doctoral work on sCO$_2$ axial-turbine blading~\citep{sathish2019sco2} used
exactly this pattern: Class-Shape-Transformation (CST) geometry, a B-spline
control-point wrapper, and a surrogate-assisted optimisation loop against
MISES. That work already used CST for blade parameterisation, together with
geometric constraints encoding expert design practice, but it used them
\emph{inside} an optimiser, and its closing section named the open problem
verbatim: integrating CST directly with a Newton-based flow solver to escape
the nested loop. The present paper's claim is that the linearity which made
those constraints convenient there is strong enough to remove the optimiser
altogether.

That escape route already exists in one solver. The modal inverse mode of
MISES \citep{drela1990viscous} appends mode-shape amplitudes directly to its
global Newton system and drives them so that computed surface speed matches a
target: a determined root-find, not a search. This is not new;
\citet{drela1990viscous} described it between 1986 and 1990. What the
built-in mode shapes of MISES are \emph{not}, however, is a complete,
geometrically meaningful, analytic design basis. They are heuristic bump
functions, not a family whose coefficients double as leading-edge radius,
trailing-edge thickness or inscribed area. Kulfan's CST
parameterisation~\citep{kulfan2008universal} is such a basis, and it is one
of the analytic families examined in the reference geometric benchmark of
\citet{masters2017geometric}. CST has also already been coupled to inverse
design once, by \citet{lane2010inverse}. There, CST is used to \emph{smooth} a
pressure-residual-driven normal-direction shape correction inside an outer
iterative loop; CST is the smoother, not an unknown in a Newton system. The
present work does not repeat either of these. It combines a property of CST
that neither prior use exploits.

\subsection{Related work}
\label{sec:related}

\paragraph{Classical inverse design.} Inverse airfoil design predates any of
the parameterisations discussed below. Lighthill's conformal-mapping
method~\citep{lighthill1945new} is the founding well-posed formulation: a
target speed distribution on the unit circle is mapped conformally to the
physical airfoil. Three integral closure conditions, closure of the contour, a
fixed leading-edge radius, and a prescribed circulation, are required for a
solvable, self-consistent shape to exist at all. This closure-condition
structure is the direct conceptual ancestor of
Section~\ref{sec:constraints}'s linear rows: both are answers to the same
question, what makes an arbitrarily chosen target realisable by \emph{some}
member of the design family, asked eighty years apart under different
geometric representations. \citet{volpe1986design} sharpened well-posedness
for the transonic regime, showing that an inverse problem posed with pressure
rather than speed as the target, and with the wrong closure conditions, is
generically ill-posed. Existence and uniqueness are not automatic once the
target is chosen independently of the geometry family, a caveat that motivates
the identifiability discussion of Section~\ref{sec:ablations} in a different,
discrete and finite-dimensional setting. Eppler and Somers's design
method~\citep{eppler1980airfoil} and Selig and Maughmer's
PROFOIL~\citep{selig1992multipoint} are the mature, still-used engineering
descendants. Both are multi-point conformal-mapping inverse solvers that
target several operating conditions' pressure distributions simultaneously by
blending mode functions, and both remain outside the Newton-coupled,
viscous-inclusive family this paper and MISES belong to. They close the
inverse problem analytically at the potential-flow level, then couple to a
boundary-layer method non-simultaneously, rather than solving flow and
geometry as one system.

\paragraph{Parameterisation comparison.} Sobieczky's PARSEC
\citep{sobieczky1999parsec} parameterises an airfoil by twelve
directly-meaningful quantities, among them leading-edge radius, the thickness
and camber extrema and their locations, and the trailing-edge angle and
thickness, fitted by a polynomial basis. It is geometrically transparent in
the same spirit as CST's endpoint identities (Eq.~\eqref{eq:endpoint}), but
through a basis whose coefficients are not independent design freedoms in the
same linear sense: PARSEC's basis functions depend on the target values
themselves, since the fit is re-solved per target, whereas CST's Bernstein
basis $S_i(\psi)$ is fixed once $n$ is chosen, independent of what shape is
being represented. This is the property Eq.~\eqref{eq:dsurface-dA} depends on.
The benchmark of \citet{masters2017geometric} is the field's reference
comparison: over two thousand airfoils and seven parameterisation families,
including CST, PARSEC, B-splines, Hicks--Henne bumps, a radial-basis
domain-element method, and orthogonal modes derived by singular-value
decomposition of a training database. Its headline finding is that twenty to
twenty-five design variables suffice for all of the families tested, with the
methods performing comparably at that budget and the data-driven modal method
marginally the most efficient. That last qualification is the one this paper
turns on. Data-driven modes fit an existing database efficiently at the cost
of needing that database, and they carry no closed-form geometric-functional
structure of the kind Section~\ref{sec:constraints} exploits. Among the
analytic families that do, CST is a well-supported choice, and it is the one
whose particular algebraic property, linearity in the coefficient vector $A$,
this paper's contribution depends on.

\paragraph{Newton-coupled viscous/inviscid solvers.} The architecture this
paper appends CST to is the ISES lineage of \citet{drela1987viscous}: a
two-equation integral boundary-layer closure coupled \emph{simultaneously},
not iteratively, to a panel-method inviscid solution through a single global
Newton system, with the $e^n$ transition model providing the closure that
makes laminar-to-turbulent transition a smooth part of that system rather than
a switched boundary condition. XFOIL~\citep{drela1989xfoil} is the widely used
low-speed member of that lineage, and \texttt{mfoil}
\citep{fidkowski2022coupled} is a modern, open, from-scratch
reimplementation of exactly this architecture: the ISES and XFOIL
viscous--inviscid coupling without XFOIL's Fortran legacy code, in Python,
with an exposed sparse global Jacobian. This is why it serves as the first
testbed here (Section~\ref{sec:formulation}): the extended system of
Section~\ref{sec:newton-system} could not be built without an analysis code
that already exposes its own $\partial R/\partial U$ block for direct
augmentation. The parametric inverse mode of MISES was, moreover, designed
from the outset to accept an arbitrary geometry-definition system through a
black-box interface \citep{drela1990viscous}, so the mechanism a later
transplant of this architecture would use already exists in the target solver,
and has never been exercised with a CST-as-unknowns formulation.

\paragraph{Modern learned inverse-design methods.} A distinct, much newer line
of work learns the pressure-to-geometry map directly from data rather than
solving a physics-based system at all. Conditional generative adversarial
networks \citep{yilmaz2020cgan} and, more recently, diffusion models learn a
stochastic generative map from a target pressure or aerodynamic-coefficient
vector to a geometry, trained offline on a large precomputed database, with
inference reduced to a single or few-step forward network pass. A
classifier-free-guided denoising diffusion model reports a 34.4 percent
precision improvement over Wasserstein-GAN baselines on an airfoil-generation
benchmark \citep{deng2025cddpm}. A hybrid approach trains a physics-informed
surrogate on low-fidelity viscous--inviscid-coupled flowfields of the same
class this paper's solver produces, and reports at least an
order-of-magnitude online inference speedup relative to direct analysis, again
after an offline training investment \citep{wong2024lfpinn}. These methods answer a different
question well: given a large library of prior designs and target statistics
resembling that library, produce a plausible geometry near-instantly. They do
not provide a convergence guarantee, an exact constraint-row Jacobian, or a
certificate that the returned geometry actually reproduces the requested
pressure distribution to any specified tolerance. The network's output is
evaluated \emph{after} the fact by running a flow solver on it, not
constrained \emph{during} generation the way Eq.~\eqref{eq:kkt}'s equality
rows are. Section~\ref{sec:comparison}'s comparison table places this paper's
architecture and the learned-inverse family on common columns explicitly,
because the two are not competing solutions to the same technical problem so
much as different points on an accuracy-guarantee against amortised-cost
trade surface.

\paragraph{Turbomachinery inverse design.} \citet{porta2026datadriven}
benchmark Kolmogorov--Arnold networks against multilayer-perceptron and
Gaussian-process regressors for inverse turbine-passage design, trained on
approximately thirty thousand blade profiles generated by an automated
MISES-coupled optimisation pipeline over inlet flow angles from $-50^\circ$
to $0^\circ$ and outlet angles from $50^\circ$ to $75^\circ$. The trained
network infers a loading-consistent geometry in about $0.1$ seconds with a
mean outlet-angle error of $0.086^\circ$. Structurally this is the same
learned-surrogate pattern as the previous paragraph, specialised to cascades,
and it is the closest published turbomachinery analogue of this paper's
motivating problem. It is, however, once more amortised-cost statistical
inference over a precomputed design database, not a per-target root-find with
a convergence certificate. Their pipeline includes a predicted-error
feasibility proxy that flags statistically ill-posed inverse targets before
training-set generation. This is the statistical-learning analogue of this
paper's realisability guard (Section~\ref{sec:presolve}), independently
motivated by the same underlying problem, that some targets are simply not
achievable by the design family in use, but answered by a learned classifier
rather than a closed-form linearisation.

\subsection{The contribution and its scope}
\label{sec:contribution}

CST's surface parameterisation
\begin{equation}
  \zeta(\psi) = C(\psi)\sum_i A_i S_i(\psi) + \psi\,\zeta_T
  \label{eq:cst-intro}
\end{equation}
is \emph{linear} in the coefficient vector $A$, and two consequences follow.
Neither is available to a generic geometry representation, such as a B-spline
with a control-point wrapper, once local-frame perturbations, leading-edge
tangent-angle displacements or re-fitting are introduced to reach the same
design freedom. First, the geometric sensitivity
$\partial\zeta/\partial A_i = C(\psi)S_i(\psi)$ is exact, constant and
design-independent, so any operator built on it is assembled once rather than
re-derived inside a Newton loop. Second, the geometric functionals that matter
for manufacturability, among them leading-edge radius, trailing-edge
thickness, wedge angle, inscribed area and leading-order curvature continuity,
are themselves \emph{linear} in $A$, so they enter the Newton system as
constant-coefficient algebraic rows rather than non-smooth predicates or
penalty terms. The contribution of this paper is not geometry unknowns in a
Newton system, which is Drela's, and not CST used somewhere inside inverse
design, which is Lane and Marshall's. It is the combination: an analytic basis
whose linearity makes the geometric sensitivity exact and constant,
\emph{and} whose geometric functionals make design constraints linear rows,
together converting constrained shape optimisation into constrained
root-finding, one square Newton system, solved once, with no outer loop.

Two scope limits on that claim are stated here rather than left for the
reader to discover. The first concerns which half of the claim the reported
experiments exercise. Every configuration reported in this paper runs either
with the leading edge prescribed, in which case no constraint row is assembled
at all, or with a single shared-leading-edge row; the trailing-edge-wedge and
inscribed-area rows are implemented and unit-tested against independent
quadrature but are not exercised by any measured result
(Section~\ref{sec:constraints}). The second concerns the Jacobian. Eq.~
\eqref{eq:dsurface-dA}'s exact constant sensitivity is used directly in the
constraint rows and in the linear pre-solve of Section~\ref{sec:presolve}; the
flow block $\partial R/\partial A$, by contrast, is obtained by finite
differences over $A$, because the vendor solver exposes no complete geometric
sensitivity of its own residual (Section~\ref{sec:derivatives}). Both limits
are consequences of what an unmodified third-party solver makes available, and
both are revisited in Section~\ref{sec:limitations}.

\subsection{The corrected headline}
\label{sec:headline}

A priori, the architecture was hypothesised to require two to three orders of
magnitude fewer flow solves, plus determinism and uniqueness, relative to a
nested-optimisation baseline. That number was never measured against a
competent baseline, and it is not repeated here. A controlled ablation study
(Section~\ref{sec:results}) instead measures the architecture against a modern,
properly-scaled, tightly-toleranced Levenberg--Marquardt
solver~\citep{virtanen2020scipy}, and finds $3.1$ to $8.1\x$ fewer counted
flow solves depending on initialisation strategy, under two independent
fair-paired controls. Because the counting currency is itself contestable,
Section~\ref{sec:cost} reports the same two pairings in measured wall-clock
time, where the advantage narrows to $3.4$--$3.5\x$; the manuscript treats the
wall-clock figure, not the solve count, as the conservative statement of the
result. The measured ratio comes with qualitative properties no cost figure
captures: determinism, with no stochastic restarts; an exact analytic Jacobian
for whatever constraint rows are active, free of finite-difference noise; and
a battery of per-iteration diagnostics (Section~\ref{sec:fm}) that discriminate
among the architecture's failure modes and for which the nested baseline has
no equivalent. The paper's two central empirical results are therefore the
correction of the a priori hundred-fold claim and the discovery of a
station-selection identifiability mechanism that the original hypothesis did
not anticipate (Section~\ref{sec:ablations}).

\subsection{Structure of the paper}

Section~\ref{sec:formulation} states the CST mathematics, the constraint rows,
the extended Newton system and its degree-of-freedom accounting, the
derivative strategy the target solver's exposed interface actually forces, and
the two solver-specific closures, forced transition and the realisability
pre-solve, needed to make the Newton iteration well-posed.
Section~\ref{sec:fm} catalogues the four failure modes the architecture is
vulnerable to and the diagnostics built to discriminate between them.
Section~\ref{sec:falsifiable} describes the self-consistency test that
falsifies or confirms the hypothesis outright. Section~\ref{sec:results}
reports the ablation matrix, the identifiability finding, the
Bernstein-conditioning cliff, the two panel evaluations, the corrected cost
comparison and the observed convergence behaviour.
Section~\ref{sec:artifacts} describes the released software artifacts, an
interactive application and a project site, through which the method can be
exercised without installing anything. Section~\ref{sec:mitigations} draws out
the common structure of the four solver-specific mitigations,
Section~\ref{sec:validity} states the study's threats to validity and
Section~\ref{sec:limitations} its scope limits, and
Section~\ref{sec:outdated-baseline} discusses why the pre-registered cost
threshold was missed. Section~\ref{sec:conclusions} closes with the cascade
and MISES extensions this architecture was designed to scale into.
